% Part: propositional-logic
% Chapter: syntax-and-semantics
% Section: valuations-sat

\documentclass[../../../include/open-logic-section]{subfiles}

\begin{document}

\olfileid{pl}{syn}{val}

\olsection{\usetoken{P}{valuation} en vervulling}

\begin{defn}[!!^{valuation}s]
De verzameling $\{\True, \False\}$ bestaat uit de twee
waarheidswaarden ``waar'' en ``onwaar''. \emph{!!^a{valuation}} voor
$\Lang{L_0}$ is een functie~$\pAssign{v}$. Deze functie kent aan elke
!!{propositional variable} van de taal $\True$ of $\False$ toe. In
symbolen: $\pAssign{v} \colon \PVar \to \{\True, \False \}$.
\end{defn}

\begin{defn}
  Neem !!a{valuation}~$\pAssign{v}$. We definiëren de evaluatiefunctie
  $\pValue{v} \colon \Frm[L_0] \to \{\True, \False \}$ stap voor stap:
  \begin{align*}
    \iftag{prvFalse}{\pValue{v}(\lfalse) & = \False; \\}{}
    \iftag{prvTrue}{\pValue{v}(\ltrue) & = \True; \\}{}
    \pValue{v}(\Obj p_n) & = \pAssign{v}(\Obj p_n); \\
    \iftag{prvNot}{\pValue{v}(\lnot !A) & = \begin{cases}
        \True & \text{als } \pValue{v}(!A) = \False;\\
        \False & \text{anders.}
      \end{cases} \\ }{}
    \iftag{prvAnd}{\pValue{v}(!A \land !B) & = \begin{cases}
        \True &
        \text{als $\pValue{v}(!A) = \True$ en $\pValue{v}(!B) = \True$;}\\
        \False &
        \text{als $\pValue{v}(!A) = \False$ of $\pValue{v}(!B) = \False$}.
      \end{cases}\\}{}
    \iftag{prvOr}{\pValue{v}(!A \lor !B) & = \begin{cases}
        \True &
        \text{als $\pValue{v}(!A) = \True$ of $\pValue{v}(!B) = \True$;}\\
        \False &
        \text{als $\pValue{v}(!A) = \False$ en $\pValue{v}(!B) = \False$}.
      \end{cases}\\}{}
    \iftag{prvIf}{\pValue{v}(!A \lif !B) & = \begin{cases}
        \True &
        \text{als $\pValue{v}(!A) = \False$ of $\pValue{v}(!B) = \True$;}\\
        \False &
        \text{als $\pValue{v}(!A) = \True$ en $\pValue{v}(!B) = \False$}.
      \end{cases}\\}{}
    \iftag{prvIff}{\pValue{v}(!A \liff !B) & = \begin{cases}
        \True &
        \text{als $\pValue{v}(!A) = \pValue{v}(!B)$;}\\
        \False &
        \text{als $\pValue{v}(!A) \neq \pValue{v}(!B)$}.
      \end{cases}}{}
\end{align*}
\end{defn}

\begin{explain}
Deze regels leveren de volgende tabellen met waarheidswaarden op:
\begin{center}
\iftag{prvNot}{
  \begin{tabular}{|c||c|} \hline
    $!A$ & $ \lnot !A$ \\
    \hline \hline
    $\True$ & $\False$ \\
    $\False$ & $\True$ \\
    \hline
  \end{tabular}
}{}
\iftag{prvAnd}{
  \begin{tabular}{|cc||c|} \hline
    $!A$ & $!B$ & $!A \land !B$ \\
    \hline \hline
    $\True$ & $\True$ & $\True$ \\
    $\True$ & $\False$ & $\False$ \\
    $\False$ & $\True$ & $\False$ \\
    $\False$ & $\False$ & $\False$ \\
    \hline
  \end{tabular}
}{}
\iftag{prvOr}{
  \begin{tabular}{|cc||c|} \hline
    $!A$ & $!B$ & $!A \lor !B$ \\
    \hline \hline
    $\True$ & $\True$ & $\True$ \\
    $\True$ & $\False$ & $\True$ \\
    $\False$ & $\True$ & $\True$ \\
    $\False$ & $\False$ & $\False$ \\
    \hline
  \end{tabular}
}{}%
\iftag{notprvNot,notprvAnd,notprvOr,notprvIf}{}{\\[1em]}
\iftag{prvIf}{
  \begin{tabular}{|cc||c|} \hline
    $!A$ & $!B$ & $!A \lif !B$ \\
    \hline \hline
    $\True$ & $\True$ & $\True$ \\
    $\True$ & $\False$ & $\False$ \\
    $\False$ & $\True$ & $\True$ \\
    $\False$ & $\False$ & $\True$ \\
    \hline
  \end{tabular}
}{}
\iftag{prvIff}{
  \begin{tabular}{|cc||c|} \hline
    $!A$ & $!B$ & $!A \liff !B$ \\
    \hline \hline
    $\True$ & $\True$ & $\True$ \\
    $\True$ & $\False$ & $\False$ \\
    $\False$ & $\True$ & $\False$ \\
    $\False$ & $\False$ & $\True$ \\
    \hline
  \end{tabular}
}{}
\end{center}
\end{explain}

\begin{prob}
Voeg aan $\Lang{L_0}$ een verbindingsteken $\diamondsuit$ voor drie
plaatsen toe. De waarde ervan wordt als volgt bepaald:
\begin{gather*}
  \pValue{v}(\diamondsuit( !A , !B , !C ) ) = \begin{cases}
    \pValue{v}( !B ) &
    \text{als $\pValue{v}(!A) = \True$;}\\
    \pValue{v}( !C ) &
    \text{als $\pValue{v}(!A) = \False $}.
  \end{cases}
\end{gather*}
Schrijf de tabel met waarheidswaarden voor dit verbindingsteken op.
\end{prob}

\begin{thm}[Lokale bepaaldheid]
\ollabel{thm:LocalDetermination} Neem twee !!{valuation}s
$\pAssign{v_1}$ en $\pAssign{v_2}$. Stel dat ze dezelfde waarde geven
aan elke !!{propositional variable} die in $!A$ voorkomt. Dus
$\pAssign{v_1}(\Obj p_n) = \pAssign{v_2}(\Obj p_n)$ telkens als
$\Obj p_n$ in de !!{formula}~$!A$ voorkomt. Dan geven ook
$\pValue{v_1}$ en $\pValue{v_2}$ dezelfde waarde aan~$!A$:
$\pValue{v_1}(!A) = \pValue{v_2}(!A)$.
\end{thm}

\begin{proof}
Door inductie naar $!A$.
\end{proof}

\begin{defn}[Vervulling]
\ollabel{defn:satisfaction} Nu definiëren we stap voor stap wanneer
  !!a{formula}~$!A$ door !!a{valuation}~$\pAssign{v}$ wordt \emph{vervuld}. We
  schrijven dit als $\pSat{v}{!A}$. (De notatie $\pSat/{v}{!A}$
  betekent ``niet $\pSat{v}{!A}$''.)
\begin{enumerate}
\tagitem{prvFalse}{%
  \indcase{!A}{\lfalse}{$\pSat/{v}{\indfrm}$.}}{}

\tagitem{prvTrue}{%
  \indcase{!A}{\ltrue}{$\pSat{v}{\indfrm}$.}}{}

\item \indcase{!A}{\Obj p_i}{$\pSat{v}{\indfrm}$
  precies als $\pAssign{v}(\Obj p_i) = \True$.}

\tagitem{prvNot}{%
  \indcase{!A}{\lnot !B}{$\pSat{v}{\indfrm}$ precies als
    $\pSat/{v}{!B}$.}}{}

\tagitem{prvAnd}{%
  \indcase{!A}{(!B \land !C)}{$\pSat{v}{\indfrm}$ precies als
    $\pSat{v}{!B}$ en $\pSat{v}{!C}$.}}{}

\tagitem{prvOr}{%
  \indcase{!A}{(!B \lor !C)}{$\pSat{v}{\indfrm}$ precies als
    $\pSat{v}{!B}$ of $\pSat{v}{!C}$ (of allebei).}}{}

\tagitem{prvIf}{%
  \indcase{!A}{(!B \lif !C)}{$\pSat{v}{\indfrm}$ precies als
    $\pSat/{v}{!B}$ of $\pSat{v}{!C}$ (of allebei).}}{}

\tagitem{prvIff}{%
  \indcase{!A}{(!B \liff !C)}{$\pSat{v}{\indfrm}$ precies als
    ofwel zowel $\pSat{v}{!B}$ als $\pSat{v}{!C}$ geldt, ofwel noch
    $\pSat{v}{!B}$ noch $\pSat{v}{!C}$ geldt.}}{}
\end{enumerate}
Als $\Gamma$ een verzameling !!{formula}s is, geldt
$\pSat{v}{\Gamma}$ precies als $\pSat{v}{!A}$ voor
elke~$!A \in \Gamma$.
\end{defn}

\begin{prop}\ollabel{prop:sat-value}
  $\pSat{v}{!A}$ precies als $\pValue{v}(!A) = \True$.
\end{prop}

\begin{proof}
  Door inductie naar~$!A$.
\end{proof}

\begin{prob}
  Bewijs \olref[pl][syn][val]{prop:sat-value}
\end{prob}
\end{document}
